\chapter{Mười Lát Cắt Người Lớn Bị Nuốt}
\markboth{Mười Lát Cắt Người Lớn Bị Nuốt}{Mười Lát Cắt Người Lớn Bị Nuốt}

\latcat{1. Cấm Con Nhưng Không Cấm Mình}{Người lớn mắng con dùng điện thoại quá nhiều rồi ngay sau đó tự kéo ghế ngồi lướt tiếp.\\
\textit{(An adult scolds a child for excessive phone use and then immediately sits down to keep scrolling.)}\\
Kiểu dạy đó không làm trẻ phục. Nó chỉ làm trẻ học cách coi thường lời nói rỗng.\\
\textit{(That style of teaching does not earn respect. It teaches children to disregard hollow words.)}}{Hypocrisy is louder than advice.}

\latcat{2. Cuộc Họp Gia Đình Bị Cắt Bởi Tin Nhắn}{Có những chuyện rất thật đang được nói trong nhà thì màn hình rung lên.\\
\textit{(Something deeply real is being discussed at home when the screen vibrates.)}\\
Và thế là cảm xúc con người thua một thông báo.\\
\textit{(And just like that, human feeling loses to a notification.)}}{Every interruption teaches someone their pain ranked second.}

\latcat{3. Dắt Con Đi Chơi Nhưng Không Thật Sự Đi}{Nhiều cha mẹ đưa con ra công viên nhưng ánh mắt lại ở trong điện thoại.\\
\textit{(Many parents take their children to the park while their eyes remain inside the phone.)}\\
Đứa trẻ có thể được đưa đi, nhưng không thật sự được ở cùng.\\
\textit{(The child may be taken out, but not truly accompanied.)}}{Presence cannot be outsourced by transportation.}

\latcat{4. Giáo Viên Giảng Về Tập Trung Nhưng Điện Thoại Luôn Bên Cạnh}{Một bài học về kỷ luật sẽ mất lực rất nhanh nếu người nói cũng rung theo màn hình.\\
\textit{(A lesson on discipline loses force quickly if the speaker vibrates with the screen too.)}\\
Học sinh nhìn rất kỹ những mâu thuẫn nhỏ đó.\\
\textit{(Students notice those small contradictions very closely.)}}{Authority leaks through inconsistency.}

\latcat{5. Trả Lời Công Việc Mãi Không Xong}{Nhiều người lớn tự nhận mình quá bận.\\
\textit{(Many adults say they are extremely busy.)}\\
Nhưng một phần lớn của cái bận đó là bận vì bị màn hình giật qua giật lại chứ không phải bận vì làm việc sâu.\\
\textit{(But much of that busyness comes from being jerked around by screens, not from doing deep work.)}}{Fragmentation often disguises itself as productivity.}

\latcat{6. Chạy Xe Cùng Điện Thoại}{Người lớn vừa lái xe vừa liếc máy rồi vẫn bảo con phải biết an toàn.\\
\textit{(Adults glance at phones while driving and still tell children to value safety.)}\\
Đó không chỉ là nguy hiểm. Đó còn là một bài học đạo đức rất méo.\\
\textit{(That is not only dangerous. It is a twisted moral lesson as well.)}}{A bad example in motion travels far.}

\latcat{7. Tập Thể Dục Nhưng Đầu Óc Vẫn Không Rời Được Máy}{Nhiều người lớn đi bộ với điện thoại, tập máy với điện thoại, nghỉ giữa hiệp cũng với điện thoại.\\
\textit{(Many adults walk with phones, exercise with phones, and even rest between sets with phones.)}\\
Họ thay đổi tư thế, nhưng không hề thay đổi trạng thái bị kéo căng liên tục.\\
\textit{(They change posture, but never change the state of constant inner tension.)}}{Movement is not the same as release.}

\latcat{8. Ngủ Cạnh Người Thân Nhưng Sống Cạnh Màn Hình}{Có những đôi vợ chồng, cha mẹ, con cái kết thúc ngày bằng hai ánh sáng lạnh nằm song song.\\
\textit{(Some spouses, parents, and children end the day with two cold lights glowing side by side.)}\\
Họ ngủ cạnh người thân nhưng thức cùng thuật toán.\\
\textit{(They sleep beside loved ones but remain awake with algorithms.)}}{Proximity is not intimacy.}

\latcat{9. Mỗi Lúc Rảnh Đều Bị Lấp Kín}{Người lớn cũng có những phút đứng chờ, ngồi đợi, đi thang máy.\\
\textit{(Adults also have moments of waiting, sitting still, riding elevators.)}\\
Nếu mọi khoảng nhỏ đều bị lấp bằng lướt, đầu óc sẽ chẳng còn chỗ thở để trưởng thành nữa.\\
\textit{(If every small gap is filled with scrolling, the mind is left with no breathing space in which to mature.)}}{No silence, no digestion of life.}

\latcat{10. Người Lớn Cũng Đang Xin Tha Bằng Ưu Điểm Nhỏ}{Nhiều người lớn nói: điện thoại giúp công việc của tôi mà.\\
\textit{(Many adults say: but my phone helps my work.)}\\
Có thể đúng. Nhưng nếu đổi lại bằng sự hiện diện rách nát, ngủ không yên, bữa cơm nguội, và những cuộc đối thoại què quặt, thì cái giá đó không còn nhỏ nữa.\\
\textit{(That may be true. But if the exchange price is shredded presence, restless sleep, cold meals, and crippled conversations, then the cost is no longer small.)}}{Useful tools become expensive when paid for with the soul of daily life.}

\ngancach